\scanpage{45}
\uline{Structure à étudier} Partition de $S$ en des $(S_\alpha)_{\alpha\in\unsure{J|I}}$ mutuellement non panachées. (On les appelle des partitions non \unsure{panachées|panachés})

% slanted
\margin{\uline{Ex} \unsure{partition|partitions} \unsure{grossière|grossières}, \unsure{ayant} \unsure{juste} une classe, égale à $S$\unsure{\ldots|--}}

\uline{Intuition} À \deleted{\unsure{une tel|un tel}} un ens. bicircularisé avec partition non panachée, \inserted{de \deleted{\unsure{Card.}} \unsure{cardinal|Cardinal} $N$,} correspond (\unsure{fonctoriellement} pour les \unsure{iso}) un ens. bicircularisé (d'\deleted{ordre} cardinal \inserted{\uline{pair}} $N'\leq N$), avec une partition \inserted{non panachée} de degré 2 (classes de cardinal 2).\unsure{\,--|\ldots}

\figure{\begin{tikzpicture}[x=.011cm,y=-.011cm,line width=.65pt]
\draw (578,1785) .. controls (632,1751) and (637,1702) .. (664,1620)
 .. controls (685,1547) and (706,1482) .. (670,1401)
 .. controls (632,1315) and (543,1273) .. (449,1281)
 .. controls (366,1284) and (291,1314) .. (225,1374)
 .. controls (173,1424) and (162,1475) .. (158,1544)
 .. controls (144,1612) and (168,1680) .. (195,1725)
 .. controls (254,1806) and (355,1865) .. (428,1858)
 .. controls (487,1855) and (556,1837) .. (601,1781);
\end{tikzpicture}}
% archive sheet number: 44
